% ============================================================================
%  Subdocumento 8. Contenido. Se usa igual en el documento único y en el
%  PDF propio del subdocumento. Los formularios que le asigna el T-21 se
%  agregan solos al final (datos en formularios/ de esta carpeta).
% ============================================================================

\begin{resumenApertura}
\marcador{Qué resuelve este subdocumento, en cuatro o cinco líneas}
\begin{recibe}
  \item \marcador{Qué recibe el mandante}
\end{recibe}
\end{resumenApertura}

\section{Identificación y cuantificación}\label{sec:8-identificacion-y-cuantificacion}

\marcador{Riesgos técnicos, organizacionales, de proyecto, de seguridad y de operación}

\section{Análisis cualitativo y cuantitativo}\label{sec:8-analisis-cualitativo-y-cuantitat}

\marcador{Modos de falla, árbol de fallas o simulación}

\subsection{Línea base acordada para capacitación y adhesión}\label{sec:8-acuerdo-capacitacion-adhesion}

La capacitación A13 conserva 40 días hábiles y usa la precedencia A06,
A09 FS$-10$ y G13 del T-15: inicio temprano en el día 182 y término en el 222.
G13 exige versión de entrenamiento probada y congelada en PREPROD, datos ficticios
y material por perfil. Las funciones certificadas al cierre del sprint 7 deben
ser identificadas antes de liberar ese hito. Los diez primeros días cubren
familiarización; los treinta restantes cubren práctica completa, evaluaciones y
complemento del sprint 8 antes de la marcha blanca. Los turnos y relevos protegen
la operación. La falta de disponibilidad de G13 retrasa A13 y debe registrarse
como atraso y actualizarse en el análisis de riesgo. La holgura total de A13 es
38 días y su holgura libre, 18; A06 tiene 24 días de holgura libre. El término
determinista de la red permanece en el día 420.

La simulación debe utilizar esta misma red y comprobar la disponibilidad de G13,
sin tratar la liberación de entrenamiento como cierta por el solo adelanto.
\marcador{Incorporar y verificar código, distribuciones, correlaciones, número de corridas,
semilla y tratamiento del atraso de G13 de la simulación de riesgos; resultados no verificados en este checkout}.

La adhesión firmada usa como denominador los 148 transportistas de la población
inicial. Se cuenta una vez por RUT, con anexo firmado por el mandante y el
representante autorizado, vigente, no revocado y registrado al corte. Se incluyen
las modalidades completa y de datos. Al cierre de M9, 31 de octubre de 2027 bajo
SA-01, se requieren al menos 60 firmados: con 59 o menos se activa el refuerzo.
Las metas posteriores son 104 en M16 y 134 en M21. Instalación, activación técnica
y consentimientos de conductores se controlan por separado; la firma del
transportista no sustituye esos consentimientos. El Formulario T-16 registra
estas condiciones sin atribuirles una cuantificación de riesgo aún no verificada.

\section{Estrategias de mitigación}\label{sec:8-estrategias-de-mitigacion}

\marcador{Costo beneficio, responsable, plazo y disparador}

\section{Riesgos tecnológicos y de contraparte}\label{sec:8-riesgos-tecnologicos-y-de-contra}

\marcador{Obsolescencia, bloqueo por proveedor, escalabilidad, ciberseguridad y disponibilidad de contrapartes}

\section{Reservas de contingencia y de gestión}\label{sec:8-reservas-de-contingencia-y-de-ge}

\marcador{Reflejo en cronograma y flujo de caja, sin montos}
